\subsection{The tail condition and perturbations of the density}

Vanishing tails are equivalent to continuity of the threshold under uniformly bounded multiplicative perturbations that converge in measure. This characterizes the precise role of the tail condition: if it fails, a perturbation supported on a shrinking set can lower the threshold by a fixed factor.

\begin{proposition}\label{prop:root-tail-continuity}
Let $0\leq p\in\mathcal D_\omega$ be nonzero. The following are equivalent:
\begin{enumerate}
\item $p\in\mathcal D_\omega^0$;
\item the densities $p k_T$, for $T\in\mathscr G_\omega$, are uniformly integrable;
\item whenever $F_j$ are smooth real functions, uniformly bounded and tending to zero in measure, one has
\[
 b_\omega(\mathrm{e}^{F_j}p^n)\longrightarrow b_\omega(p^n).
\]
\end{enumerate}
If the first condition holds, the convergence in the third also holds for bounded measurable $F_j$ with the same uniform bound and convergence in measure. If it fails, there are smooth $F_j\geq0$, uniformly bounded and tending to zero in measure, such that
\[
 \mathrm{e}^{F_j}p^n\longrightarrow p^n\quad\text{in }L^1,
 \qquad
 b_\omega(\mathrm{e}^{F_j}p^n)\leq2^{-n}b_\omega(p^n).
\]
\end{proposition}
\begin{proof}
The first condition implies the second by \cref{thm:integrable-current-attainment}. Conversely, $\mathrm{d}V(\{p>K\})\to0$, so uniform integrability and \eqref{eq:integrable-current-supremum} give
\[
 \rho_\omega(p\mathbf1_{\{p>K\}})
 =\sup_{T\in\mathscr G_\omega}\int_{\{p>K\}}p k_T\,\mathrm{d}V
 \longrightarrow0.
\]
This is the first condition by \cref{prop:root-uniform-tails}.

If $p\in\mathcal D_\omega^0$ and $|F_j|\leq A$, then $p(\mathrm{e}^{F_j/n}-1)$ tends to zero in measure and its absolute value is at most $(\mathrm{e}^{A/n}+1)p$. The dominated-convergence criterion of \cref{prop:root-uniform-tails} makes its root norm tend to zero. The triangle inequality gives $\rho_\omega(p\mathrm{e}^{F_j/n})\to\rho_\omega(p)>0$. Taking negative $n$th powers proves the third condition, including measurable perturbations.

Suppose now that the tails do not tend to zero, and put $R=\rho_\omega(p)>0$. The smooth Gauduchon formula yields $\delta>0$, numbers $K_j\to\infty$ and smooth strictly positive $T_j\in\mathscr G_\omega$ such that
\[
 \int_{\{p>K_j\}}p k_{T_j}\,\mathrm{d}V>\delta.
\]
The finite measure $p k_{T_j}\,\mathrm{d}V$ is regular. Choose a compact set $C_j\subset\{p>K_j\}$ whose mass exceeds $\delta/2$, and an open neighborhood $O_j$ with
$\mathrm{d}V(O_j)\leq\mathrm{d}V(\{p>K_j\})+1/j$. Take a smooth cutoff $0\leq\chi_j\leq1$ that equals one on $C_j$ and is supported in $O_j$. Choose a fixed $A>0$ with $\mathrm{e}^{A/n}\delta/2>2R$, and set $F_j=A\chi_j$. Then $F_j\to0$ in measure, while
\[
 \rho_\omega(p\mathrm{e}^{F_j/n})
 \geq\int_Xp\mathrm{e}^{F_j/n}k_{T_j}\,\mathrm{d}V
 >2R.
\]
The threshold formula gives the asserted uniform gap. Finally,
$|\mathrm{e}^{F_j}-1|p^n\leq(\mathrm{e}^A-1)p^n\mathbf1_{O_j}$, whose integral tends to zero by absolute continuity of the integral. This proves the last assertion and the converse to the third condition.
\end{proof}

For $f=p^n$, these perturbations satisfy $\mathrm{e}^{-A}f\leq\mathrm{e}^{F_j}f\leq\mathrm{e}^Af$ with a fixed $A$. The continuity statement is restricted to such an order interval. The unrestricted $L^1$ instability in \cref{sec:individual-density-stability} allows approximations without these bounds.

For attainment at one density, uniform integrability of a single maximizing sequence is enough. This weaker condition also identifies the limit of its weighted measures.

\begin{proposition}\label{prop:selected-gauduchon-attainment}
Let $0\leq p\in\mathcal D_\omega$ be nonzero and put $R=\rho_\omega(p)$. Suppose $T_j\in\mathscr G_\omega$ converge weakly to $T$ and satisfy $D_p(T_j)\to R$. If the densities $p k_{T_j}$ are uniformly integrable, then
\[
 D_p(T)=R,\qquad
 p k_{T_j}\,\mathrm{d}V\rightharpoonup p k_T\,\mathrm{d}V.
\]
Consequently, a uniformly integrable maximizing sequence has a maximizing current among its weak cluster limits.
\end{proposition}
\begin{proof}
For positive integers $K$, put $p_K=\min(p,K)$. Since $p$ is finite almost everywhere, uniform integrability gives
\[
 \varepsilon_K:=\sup_j\int_{\{p>K\}}p k_{T_j}\,\mathrm{d}V
       \longrightarrow0.
\]
For every nonnegative continuous $h$, the weight $hp_K$ is bounded and measurable. The bounded-weight upper semicontinuity proved in \cref{thm:integrable-current-attainment} therefore gives
\[
 \limsup_j\int_Xhp k_{T_j}\,\mathrm{d}V
 \leq \int_Xhp_K k_T\,\mathrm{d}V
                   +\|h\|_\infty\varepsilon_K.
\]
Let $K\to\infty$. For $h=1$ this yields $R\leq D_p(T)$; the supremum formula \eqref{eq:integrable-current-supremum} gives the reverse inequality. If $\mu$ is a weak cluster limit of $p k_{T_j}\,\mathrm{d}V$, the same estimate for all nonnegative continuous $h$ gives $\mu\leq p k_T\,\mathrm{d}V$. Both measures have mass $R$, so they agree. Their bounded masses give weak compactness, and every cluster limit has now been identified. The final assertion follows from compactness of $\mathscr G_\omega$.
\end{proof}
